\begin{lemma}[Three-dimensional quotient algebras]\label{B:small-quotient-algebra}
Let $A$ be a unital finite-dimensional $\mathbb Q$-algebra with
$\dim_{\mathbb Q}A\leq3$, let $W$ be a finite-dimensional right
$A$-module, and let $V$ be a finite-dimensional left $A$-module.
If $\operatorname{Tor}_2^A(W,V)\ne0$, then
\[
  \operatorname{Tor}_n^A(W,V)\ne0\qquad\text{for every }n\geq2.
\]
Consequently, if these Tor groups vanish in all sufficiently high
degrees, then $\operatorname{Tor}_2^A(W,V)=0$.
\end{lemma}

\begin{proof}
Write $J=\operatorname{rad}(A)$.  If $J=0$, the algebra is semisimple
and all positive Tor groups vanish.  Suppose $J\ne0$.
The dimension bound gives the following complete list of possibilities:
\[
 \mathbb Q[x]/(x^2),\quad
 \mathbb Q\times\mathbb Q[x]/(x^2),\quad
 \mathbb Q[x]/(x^3),\quad
 \mathbb Q\oplus J\ (\dim J=2, J^2=0),\quad
 T_2(\mathbb Q).
 \tag{\ref{B:small-quotient-algebra}.1}
\]
Here $T_2(\mathbb Q)$ denotes the algebra of upper triangular
$2$ by $2$ matrices.  We first justify the list, including the possible
noncommutative case.

If $A/J\cong\mathbb Q$, the scalar copy of $\mathbb Q$ splits the
quotient and $A=\mathbb Q\oplus J$ as a vector space.  A
one-dimensional nilpotent radical has square zero and gives the dual
numbers.  For a two-dimensional radical, either $J^2=0$, or
$\dim J^2=1$ and $J^3=0$.  In the latter case the nonzero
multiplication
$(J/J^2)\otimes(J/J^2)\longrightarrow J^2$ implies that a lift
$x$ of any nonzero element of $J/J^2$ has $x^2\ne0$.
Thus $1,x,x^2$ is a basis and $x^3=0$, giving
$\mathbb Q[x]/(x^3)$.

The only remaining nonsemisimple case has
$\dim(A/J)=2$ and $\dim J=1$.  The semisimple quotient cannot be
a quadratic field: the nonzero radical would then be a one-dimensional
$\mathbb Q$-vector space carrying a unital action of that field.
Hence $A/J\cong\mathbb Q\times\mathbb Q$ and $J^2=0$.
Lift the two orthogonal idempotents of the quotient to $e$ and $1-e$.
For completeness, if $p$ is any lift of the first idempotent and
$\delta=p^2-p\in J$, then
$e=p+(1-2p)\delta$ is idempotent because $J^2=0$.
The one-dimensional bimodule $J$ has exactly one nonzero component
among $eJ e$, $eJ(1-e)$, $(1-e)Je$, and $(1-e)J(1-e)$.
A diagonal component gives a product of the dual numbers with
$\mathbb Q$; an off-diagonal component gives $T_2(\mathbb Q)$,
up to interchanging the two idempotents.  This proves the list.

The triangular algebra has global dimension at most one.  One can see
this without a classification theorem: its modules are representations
of a quiver consisting of one arrow.  For a representation
$f:V_1\longrightarrow V_2$, let $P_1=(\mathbb Q\to\mathbb Q)$
and $P_2=(0\to\mathbb Q)$ be the two vertex projectives.  The
sequence
\[
 0\longrightarrow P_2\otimes V_1
 \longrightarrow (P_1\otimes V_1)\oplus(P_2\otimes V_2)
 \longrightarrow (V_1\xrightarrow{f}V_2)\longrightarrow0
\]
is exact; at the second vertex its maps are
$v\mapsto(v,-f(v))$ and $(v,w)\mapsto f(v)+w$.
Thus this case cannot have nonzero $\operatorname{Tor}_2$.

Consider next $A=\mathbb Q[x]/(x^m)$ with $m=2$ or $3$.
Every finite-dimensional module is a direct sum of the cyclic modules
$A/(x^a)$, $1\leq a\leq m$, by the Jordan decomposition of its
nilpotent $x$-action.  The summand with $a=m$ is free.  If
$1\leq a<m$, the module $A/(x^a)$ has a free resolution whose
successive differentials alternate between multiplication by
$x^a$ and $x^{m-a}$.  Tensoring with $A/(x^b)$, for
$1\leq b<m$, gives in every positive degree homology of dimension
\[
  \min(a,b)+\min(m-a,b)-b
   =\min\{a,b,m-a,m-b\}>0.
\]
Positive Tor therefore vanishes when either cyclic summand is free,
and is nonzero in every positive degree when both are nonfree.
Direct sums prove the assertion for these algebras.  They also prove
it for $\mathbb Q\times\mathbb Q[x]/(x^2)$, since Tor splits by
the two central idempotents and the $\mathbb Q$ factor contributes
no positive-degree Tor.

Finally suppose $A=\mathbb Q\oplus J$, $\dim J=2$, and $J^2=0$.
Set $k=A/J$.  A minimal free cover of a finite module has kernel
contained in the radical of the free module.  Thus, for a nonprojective
left module $V$, its first syzygy is a nonzero semisimple module
$\Omega V\cong k^r$, $r>0$.  The same holds on the right.
In particular, if a right module $W$ is nonprojective, its minimal
free resolution never terminates: the first syzygy is nonzero
semisimple, and the syzygy of $k^s$ is $J^s\cong k^{2s}$.
All differentials in a minimal resolution become zero after
tensoring with $k$, so
$\operatorname{Tor}_i^A(W,k)\ne0$ for every $i\geq1$.
The hypothesis $\operatorname{Tor}_2^A(W,V)\ne0$ forces both
$W$ and $V$ to be nonprojective.  Dimension shifting now gives
\[
 \operatorname{Tor}_n^A(W,V)
 \cong\operatorname{Tor}_{n-1}^A(W,\Omega V)
 \cong\operatorname{Tor}_{n-1}^A(W,k)^r\ne0
 \qquad(n\geq2).
\]
This settles the last case and proves the lemma.
\end{proof}

\begin{corollary}\label{B:four-outside-minimal-band}
Let $M$ be a finite monoid whose minimal ideal $K$ is a rectangular
band.  If $BM$ has finite-dimensional homotopy type and
$H_4(BM;\mathbb Q)\ne0$, then $|M\setminus K|\geq4$.
\end{corollary}

\begin{proof}
Put $A=\mathbb QM/\mathbb QK$, so
$\dim_{\mathbb Q}A=|M\setminus K|$.
The natural dimension-shifting identification in
Lemma~\ref{lem:B-band-Tor} expresses $H_n(BM;\mathbb Q)$, for
$n\geq2$, as $\operatorname{Tor}_{n-2}^A(W,V)$ for the finite
row and column augmentation modules $V,W$.
If $\dim A\leq3$, nonzero $H_4$ and
Lemma~\ref{B:small-quotient-algebra} would force nonzero
$H_n(BM;\mathbb Q)$ for every $n\geq4$, contradicting finite
homotopy dimension.
\end{proof}
